% =====================================================================
%  08_reentrenamiento.tex — Reentrenamiento eficiente de los modelos
%  Fuente técnica: rama mejora/reentrenamiento-eficiente
%  (paquete reentrenamiento/ y carpeta docs/ del repositorio).
% =====================================================================
\chapter{Reentrenamiento eficiente de los modelos}
\label{cap:reentrenamiento}

Este capítulo describe cómo actualizar los tres modelos de \saru{}
(M1 autobús, M2 zona del número y M3 dígitos) con la menor revisión
manual posible y sin poner en riesgo la operación del patio. Lo que aquí
se marca \emph{Por validar} está implementado y probado con datos
sintéticos, pero todavía no se ha medido con datos reales del patio.

\section{Problema}

Hasta el Prototipo~1, actualizar un modelo implicaba revisar y etiquetar
a mano todos los recortes guardados. Al 19 de noviembre de 2026 había
6\,852 recortes de la vista delantera de M3 sin etiquetar en Roboflow,
con quince o más casi idénticos por cada pasada de un camión. El criterio
anterior pedía un polígono por dígito, alrededor de un minuto por imagen,
aunque M3 es un detector de cajas; el criterio vigente se fija en la
sección~\ref{sec:criterio-etiquetado}. El repositorio no tenía código de
entrenamiento, evaluación ni pruebas, y los modelos en producción se
compilaron a TensorRT en otro equipo~\cite{tensorrt}.

\begin{aviso}
La exactitud del número completo en camiones CCA fue de 12 de 16
(75\,\%) en una muestra revisada a mano. La muestra es pequeña: el valor
sirve como referencia, no como medición formal.
\end{aviso}

\section{Proceso propuesto}

El proceso completo es: video, eventos, extracción, filtrado, selección
diversa, preanotación, revisión humana selectiva, dataset versionado,
ajuste fino (\emph{fine-tuning}), evaluación y despliegue controlado.
Cada paso es un comando de \texttt{python -m reentrenamiento}, y el
cuaderno de Google Colab solo llama a esos mismos comandos.

\figuraPendiente{09\_reentrenamiento/flujo-general.png}
  {Flujo general del reentrenamiento con los dos puntos de aprobación
   humana (exportar de \texttt{docs/diagramas/01\_flujo\_general.mmd}).}
  {fig:reentr-flujo}

\section{Selección de fotogramas clave}

Los recortes se agrupan por \emph{pasada}: misma zona, mismo
identificador de seguimiento de ByteTrack~\cite{zhang2022bytetrack} y
menos de seis segundos entre recortes. El identificador ya viene en el
nombre del archivo que guarda producción. Dentro de cada pasada se
comparan con un \emph{hash} perceptual (dHash de 64 bits) para quitar
repeticiones. Con el modelo actual se mide la incertidumbre de cada
recorte, siguiendo la idea del aprendizaje activo~\cite{settles2009}.

\begin{longtable}{@{}L{3.6cm}L{10.4cm}@{}}
  \caption{Razones para mandar un recorte a revisión humana.}
  \label{tab:reentr-razones} \\
  \toprule
  \textbf{Razón} & \textbf{Regla} \\
  \midrule
  \endfirsthead
  \toprule
  \textbf{Razón} & \textbf{Regla} \\
  \midrule
  \endhead
  \texttt{sin\_digitos} & El modelo no detectó ningún dígito. \\
  \texttt{digitos\_raros} & Detectó una cantidad distinta de 3 o 4. \\
  \texttt{desacuerdo} & Su lectura no coincide con el consenso de la pasada. \\
  \texttt{baja\_conf} & Algún dígito con confianza menor a 0.50. \\
  \texttt{sistema\_distinto} & Producción registró otro número; por ejemplo,
    corrigió \texttt{0063} a \texttt{5063}. \\
  \texttt{fuera\_padron} & La lectura no está en el padrón. \\
  \texttt{condicion\_nueva} & Luz de noche o contraluz. \\
  \bottomrule
\end{longtable}

Se descartan los recortes borrosos, oscuros, sobreexpuestos, sin
contraste o demasiado pequeños, los repetidos y los que exceden el cupo
de tres por pasada. Para no entrenar solo con casos difíciles, un 20\,\%
de las pasadas sin problemas se conserva como muestra de control.

\section{Preanotación y revisión}

El modelo actual dibuja las cajas con un umbral bajo (0.25), porque es
más rápido borrar una caja que dibujarla. El lote se sube a Roboflow en
formato YOLO, y cada imagen lleva la razón por la que fue elegida para
revisar primero los casos más útiles. Una segunda persona revisa el
10\,\% al azar; si encuentra más de 5\,\% de errores, el lote se devuelve
completo.

\section{Criterio de etiquetado}
\label{sec:criterio-etiquetado}

M3 es un detector de cajas: cada dígito del número económico lleva una
caja delimitadora con su clase, de \texttt{0} a \texttt{9}. No se usa una
caja para el número completo ni un polígono por dígito.

\begin{itemize}
  \item \textbf{Una caja por dígito.} Nunca una sola caja para dos dígitos
    ni para el número completo.
  \item \textbf{Caja ajustada.} El contorno toca el trazo del dígito con
    1--2 px de margen y no invade al vecino.
  \item \textbf{Todos los dígitos legibles.} Un dígito sin anotar enseña al
    modelo que ese trazo es fondo y produce números incompletos al armar la
    secuencia.
  \item \textbf{Solo el número económico.} No se etiquetan letras, logotipos,
    placas, teléfonos ni otros números de la carrocería.
  \item \textbf{Clases exactas.} \texttt{0} a \texttt{9}, escritas igual en
    todas las imágenes y en el mismo orden que espera
    \path{modelo_digitos_v3}.
\end{itemize}

El polígono se descartó por dos razones. Primero, al exportar a YOLO
Roboflow lo convierte en una caja alineada a los ejes, así que el ajuste a
la inclinación se pierde. Segundo, dibujarlo a mano tarda alrededor de un
minuto por imagen, frente a unos segundos con cajas preanotadas. La
inclinación por perspectiva se cubre con rotaciones pequeñas en los
aumentos, no con el contorno.

En los casos límite se sigue el mismo criterio de legibilidad: un dígito
cortado por el borde se etiqueta si se ve más de la mitad y su clase es
reconocible; uno borroso, solo si una persona lo lee con seguridad; una
imagen sin ningún dígito legible se elimina del dataset o se marca como
nula, nunca se deja con anotaciones dudosas.

\section{Dataset versionado y entrenamiento}

El dataset se divide en entrenamiento, validación y prueba \emph{por
pasada}: los recortes de un mismo camión nunca quedan en dos conjuntos.
El conjunto de prueba se \textbf{congela} la primera vez y se reutiliza
en cada versión, para que todas se comparen contra el mismo examen.

El entrenamiento se hace en Google Colab Pro con
Ultralytics~\cite{ultralytics} 8.4.171, la misma versión que la Jetson.
Se parte de los pesos actuales (v3) con una tasa de aprendizaje baja y
se combinan los datos anteriores con los nuevos. Cada corrida queda
registrada con el \emph{commit}, el \emph{hash} del dataset, las
versiones de software, los argumentos y las métricas.

\section{Evaluación y criterios de aceptación}

El modelo candidato se compara contra v3 sobre el conjunto congelado.
Se \textbf{rechaza} si la precisión, el recall o el mAP bajan más de un
punto, si cualquier clase pierde más de tres puntos de recall (olvido),
o si en M3 baja la exactitud del número completo o sube el CER. Solo se
acepta si además mejora al menos una métrica principal. Después, en la
Jetson, se prueba con video grabado (conteo, seguimiento, lectura y
FPS) antes de activarlo.

\section{Despliegue y regreso}

El \texttt{.engine} se compila en la propia Jetson en FP16, junto al
modelo actual y sin reemplazarlo. Se activa cambiando una variable del
\texttt{.env}, y regresar a la versión anterior toma menos de un minuto.
Hay dos aprobaciones humanas: antes de compilar y antes de activar.

\figuraPendiente{09\_reentrenamiento/bloques.png}
  {Diagrama de bloques: infraestructura, inferencia, reentrenamiento y
   administración (exportar de
   \texttt{docs/diagramas/06\_bloques\_arquitectura.mmd}).}
  {fig:reentr-bloques}

\section{Estado}

\begin{nota}
\textbf{Implementado:} selección, preanotación, validación, división,
registro de entrenamiento, evaluación, cuaderno de Colab y script de
exportación, con 45 pruebas automáticas del paquete y 15 que documentan
el comportamiento actual de \texttt{validar\_numero}. \textbf{Por validar:} porcentaje
real de recortes descartados, tiempo de etiquetado con cajas
pre-dibujadas, primer M3~v4 y ganancia de velocidad con FP16. Nada se ha
desplegado en la Jetson.
\end{nota}
